\section{The question: which nearby measurement is informative?}
An operating map rarely arrives as an evenly spaced grid. Some measurements
are missing, accessible operating points obey constraints, and a new
measurement may be expensive. Given coordinates $\vect x_i\in\R^d$ and scalar
observations $f_i$, we would like to estimate a nearby direction of increase
or decrease. A second question is whether that direction is accurate enough
to use. These questions require different evidence.
\glsadd{sym:sampleposition}\glsadd{sym:scalarfield}

For machine characterization, the coordinates might be normalized direct,
quadrature, and excitation currents, and $f$ might be torque, a loss, or an
efficiency-related objective. Such an objective is a scalar field; calling it
a potential does not establish a conservation law. Physical flux maps
derived from magnetic co-energy have additional structure that a generic
torque or loss map does not share.

The simplest local reconstruction asks what affine function passes through
the vertices of a simplex. Its slope is an interpretable finite difference
in several directions at once. Delaunay triangulation supplies those
neighborhoods. Interpolation creates a continuous piecewise-affine scalar
field, whose gradient is constant within each cell and jumps between cells.
The derivative values can then be represented at cell centroids and
triangulated again. This creates a new interpolant, with a new modeling
choice at each recursion.

Gradient-Iso adds a geometric diagnostic: intersect level sets of each
derivative component with edges of the derivative mesh, then measure the
proximity of those intersections to its support points. The crucial issue is
whether this proximity measures reliability, or merely the shape and sampling
of derivative variation. We derive the construction before interpreting it.

\paragraph{What is established here.}
The first-order solve reproduces affine fields, its sensitivity follows from
the simplex edge matrix, and the edge cuts satisfy a linear interpolation
identity. Recursion provides an implementable higher-order estimator, but
does not by itself give accuracy guarantees. Iso coverage has counterexamples
to a general confidence interpretation. Synthetic experiments evaluate these
statements separately; they do not validate an industrial measurement loop.

This is a research revision of the earlier MATLAB laboratory and IEEE-style
draft. Historical sources are retained outside the active paper. Python
uses the existing VICE geometry backend rather than another port of the
Delaunay search routines. The RooTri draft in this repository addresses edge
extraction; this paper concentrates on the differential reconstruction.

\paragraph{Relation to derivative recovery.}
Recovering derivatives from local interpolants belongs to a broader numerical
analysis literature. Polynomial preserving recovery provides a useful
comparison: its accuracy statements depend on explicit mesh and recovery
assumptions \cite{naga2004}. Our centroid recursion is a different operator,
and those results cannot be transferred to it. We include a local quadratic
fit as an independently implemented comparator, not as a claim that it is
identical to an established recovery scheme.

\begin{figure}[htbp]
\centering\begin{tikzpicture}[node distance=7mm]
\node[pipeline block,model block] (samples) {Scattered samples $(\vect x_i,f_i)$\\declared coordinate metric};
\node[pipeline block,below=of samples] (fit) {Delaunay cells and anchored edge solves\\$\matr E_\sigma\vect g_\sigma=\Delta\vect f_\sigma$};
\node[pipeline block,below=of fit] (cloud) {Derivative samples at representatives\\new mesh at each order};
\node[diagnostic block,below left=10mm and 0mm of cloud] (search) {Gradient direction\\candidate measurement};
\node[diagnostic block,below right=10mm and 0mm of cloud] (iso) {Derivative level cuts\\iso coverage descriptor};
\draw[flow] (samples)--(fit);
\draw[flow] (fit)--(cloud);
\draw[flow] (cloud)--(search);
\draw[flow] (cloud)--(iso);
\draw[flow,guide pattern] (cloud.east)--++(1.5,0)|-(fit.east)
 node[annotation,pos=.7] {recurse};
\end{tikzpicture}

\caption{A derivative proposes a direction; an iso descriptor measures the
proximity of selected geometric cuts. Accuracy requires additional evidence.
Each higher-order stage changes both support and interpolant.}
\label{fig:pipeline}
\end{figure}
